\section{Estrategia de pruebas}
\textbf{Propósito:} Que el equipo de desarrollo sepa cuáles son los tipos de pruebas que se van a realizar y la estrategia general para llevar a cabo esas pruebas.

\textbf{Contenido:}
\begin{itemize}
  \item Explicar los tipos de pruebas (Funcionales y no funcionales) que se van a realizar a los distintos elementos que estarán bajo pruebas. Ejemplo: Pruebas unitarias, de integración, de sistema, de regresión, de rendimiento y de aceptación del usuario. Debe incluir al menos dos tipos de pruebas en su Plan de Pruebas.
  \item Plantee la estrategia para realizar cada tipo de prueba. Especifique la(s) herramienta, ambiente, criterios de completitud y otra información que considere relevante. Incluya referencias a las herramientas utilizadas (e.g. JUnit, Selenium, Mockito, Bugzilla, etc.).
\end{itemize}

\textbf{Tamaño sugerido:} 3 páginas.